% GENERATED from draft/intro.md by analysis/md2tex.py — review before trusting.
% Verified: every \citep key below resolves in manuscript/refs.bib.

\section{Introduction}



The pharmacology of a G protein-coupled receptor is defined by its conformational state rather
than by its fold. Two ligands that bind the same pocket of the same receptor can produce opposite
physiology, and the difference is carried not by the backbone the receptor adopts on average but
by which cytoplasmic conformation it populates and for how long: an analysis of the deposited class
A structures concludes that ``biased signaling is due to ligand-specific population changes in the
activation landscape, rather than to completely distinct ligand-induced conformations''
\citep[p.~8]{paajanen2026activation}. That conformational description
has become steadily less binary as the experimental methods have improved. A recent survey of
class A receptors sets out the modern picture directly, describing receptors that occupy ``three
different conformational states (active, inactive, intermediate-active) that can be interconverted
for the activation/inactivation process according to a multistate, rheostat-like model, instead of
a binary (on/off) switch model'' \citep[p.~3696]{georgiou2025heterogeneity}, with conformational heterogeneity
present even within what is conventionally called the inactive ensemble
\citep[p.~3699]{georgiou2025heterogeneity}. The same conclusion emerges from the structural databank when it
is analysed without supervision. Projecting the deposited class A structures onto a single learned
activation coordinate shows that apo and agonist-bound receptors are each bimodal, populating both
inactive and active conformations, while transducer-bound structures fall almost entirely on the
active side \citep{paajanen2026activation}.

Which input sets that state is therefore a substantive biological question rather than a modelling
convention, and the answer is not simply the ligand. Occupancy of the orthosteric pocket and
occupancy of the active state are distinct variables whose coupling varies from receptor to
receptor, strongly in A2AR and weakly in $\beta_1$AR, $\beta_2$AR and the $\mu$-opioid receptor
\citep[p.~3698]{georgiou2025heterogeneity}. In class B GCGR the two decouple almost entirely: ``agonist binding
alone is insufficient to promote TM6 opening'', and the outward movement of TM6 appears only once
Gs engages \citep[p.~1]{hilger2020gcgr}. Read across the databank the same asymmetry holds, with ``agonist
binding shifts this conformational ensemble towards the active state but does not fully stabilize
it. Instead, a stable active state is only established upon G protein binding, which locks the
receptor in its active conformation'' \citep[p.~1]{paajanen2026activation}. The contact that carries this is
specific rather than diffuse. It is the $\alpha_5$ helix of G$\alpha$, and in particular its distal
C-terminal segment, which ``must be inserted into the receptor's cytoplasmic cleft to couple with
the GPCR'' \citep[p.~3694]{georgiou2025heterogeneity}. A receptor's state is set by what is bound at the
intracellular face at least as much as by what is bound in the pocket.

Physical simulation says the same thing in the language of free energy, and adds a detail that
matters for what follows. Activation is a transition ``that span[s] multiple scales, from local
rearrangements of side chains and hydration sites to large-scale reorganizations of the
transmembrane helices'' \citep[Introduction]{aureli2026epath}, and it can now be characterised as a
continuous free-energy landscape for apo class A receptors, using the conserved PIF, DRY, NPxxY and
YY microswitch motifs as quantitative descriptors rather than as a call by eye
\citep[Methods]{aureli2026epath}. Those landscapes are not flat. Two milliseconds of aggregate dynamics
on the $\beta_2$-adrenergic receptor, aggregated into a Markov state model, show the agonist condition
reaching active-state conformations while ``inverse agonist and apo simulations do not sample active
state conformations along all structural metrics'' \citep[Results]{kohlhoff2014gpcr}. In a physical
ensemble, the apo receptor largely does not get there on its own.

Nor is there a single order in which the two inputs arrive, and the distinction matters for what a
prediction can be asked to represent. One route is agonist-first: an analysis of the deposited
class~A structures concludes that activation is ``mainly regulated by conformational selection at
the ligand level and complemented by the induced fit mechanism at the transducer level''
\citep[p.~3]{paajanen2026activation}, with ``ligand binding\ldots{} the first event in a stimulated
activation sequence'' and the outward movement of TM6 treated as ``a prerequisite for G protein
binding'' \citep[pp.~3, 6]{paajanen2026activation}. Kinetic measurements at GCGR describe the same
ordering as two steps rather than one, where agonist binding ``promotes the initial formation of the
receptor--G protein complex and subsequent full engagement of the G protein at a slower time scale''
\citep[p.~1]{hilger2020gcgr}. The other route inverts it. Receptor and G protein can associate
before agonist arrives, and the resulting pre-coupled complex is not the fully active state but an
activation intermediate: the I1$^{\mathrm{TM6}}$ conformation ``corresponds to the complex
GPCR$-$G'', is more stable with GDP bound than without, and is separated from the nucleotide-free
ternary complex by nucleotide state rather than by agonist, with GDP removal shifting the
equilibrium onward \citep[pp.~10--12, 19]{georgiou2025heterogeneity}.

Whether that second route reflects a genuinely inactive receptor is contested, and the distinction
matters. Pre-assembly, in which a G protein is bound to a receptor that has not been activated, is
not the same as basal coupling, in which the G protein engages a receptor that has spontaneously
adopted an active conformation. Inactive-state pre-assembly has been reported for Gq-coupled receptors and, earlier, for several
receptors in living cells, on resonance-energy-transfer evidence. Applying two-photon polarization
microscopy, which
needs only a single fluorescent label and therefore avoids the controls that complicate resonance
energy transfer, Bondar and Lazar found no pre-assembly between Gi1 and four receptors; the one
apparent interaction, at CB1R, disappeared when basal activity was removed with an inverse agonist
or an inactivating mutation, and they bound any residual pre-assembled fraction at endogenous
expression to below roughly 3\% of receptor molecules \citep{bondar2017preassembly}. Whether a
receptor and transducer associate before activation is therefore unsettled, and may differ by G
protein family.

Co-folding models have no representation of either sequence. Every component is supplied at once,
so a predicted complex is order-agnostic by construction and cannot adjudicate between the routes,
still less between pre-assembly and basal coupling. That is a real limit on what any prediction
result can claim about mechanism, and it cuts both ways: a co-folding experiment cannot settle the
question, but neither is it obliged to commit to an answer. Supplying a receptor with a transducer
element and no agonist is best understood as a controlled input condition rather than as a model of
any particular physiological intermediate.

Structure prediction has meanwhile become accurate enough to be used as a working instrument in
receptor pharmacology, with co-folding models reproducing GPCR backbones to within about 1.4 \AA{} even
where their ligand poses remain poor \citep{zhang2026generalization}, and it represents none of this. Co-folding models of the AlphaFold 3 lineage
return one structure per input, a limitation their authors state plainly: such models ``typically
predict static structures as seen in the PDB, not the dynamical behaviour of biomolecular systems
in solution'', and ``this limitation persists for AF3, in which multiple random seeds for either the
diffusion head or the overall network do not produce an approximation of the solution ensemble''
\citep[p.~498]{abramson2024af3}. The failure is not merely one of breadth. On the paired apo and ligand-bound
forms of cereblon, ``AF3 exclusively predicts the closed state for both holo and apo systems''
\citep[p.~498]{abramson2024af3}, which is to say the model returns whichever state dominates its training
data regardless of what the input specifies. Independent work has since established that this is
the general behaviour rather than an anecdote. On a PDB-wide, mechanism-stratified benchmark, four
co-folding models and one ensemble emulator recover all known intrinsic states in only about 8 to
29\% of sequence clusters, and the bottleneck is traced to the structure module rather than to any
loss of diversity in the pair representation \citep{ku2026promise}. Across four multi-state proteins the
models ``default to a dominant state represented in the PDB'', and the bias survives every
architecture tested \citep[p.~2]{ye2026multistatebias}. On human protein kinases, state accuracy plateaus
near 65 to 75\%, ensembles are bimodally all-correct or all-incorrect across samples, and the usual
geometric quality scores do not predict which \citep{sun2026kinconfbench}. On SLC transporters the
deposited state is returned whatever the sampling intervention, to the point that ``the ESM-AF2
protocols described here fail to generate the alternative conformational state when one
conformational state was available in the PDB at the time of AF2 training''
\citep[p.~11]{swapna2025memorization}. The most quantitative statement of the effect comes from enzymes.
Across 82 enzymes with a known open-closed domain motion, 500 AlphaFold 3 models per condition and
no templates, whether the model returns the open or the closed form is governed by how many apo and
holo structures of that enzyme sit in the PDB rather than by whether the triggering ligand is
supplied: adding the ligand moves the closed fraction by 9.1 to 17.5 percentage points, against a
40.3-point gap between the enzymes whose databank entries are mostly apo and those whose entries are
mostly holo \citep[Discussion]{yu2026domainmotion}. The authors' own summary is that ``results may depend
more on the number of structures available for training than on the presence of a ligand in the
calculation'' \citep[Significance]{yu2026domainmotion}, and the same behaviour reproduces in AlphaFold 2.

For receptors the collapse has a direction, and it runs against the state that matters most for
agonist pharmacology. Sequence-only AF2 and AF3 models agree best with inactive references and
agree progressively less well as the reference becomes more active \citep{chib2025gpcrstates}; the
earlier statement of the same phenomenon is that AF2 ``only predicts one state and is biased toward
either the active or inactive conformation depending on the GPCR class'' \citep[p.~1873]{heo2022multistate}.
Nor is this a problem that could be solved by generating many structures and selecting among them,
because the models' own confidence does not track conformational correctness. Cfold reports ``no
clear relationship between the plDDT and known conformations from the PDB suggesting that
confidence metrics can't be used to select for certain conformations'' \citep[p.~5]{bryant2024cfold}; the
multi-state survey concludes that pLDDT and PAE ``are not reliable selectors of biologically or
physically meaningful alternative states'' \citep[p.~15]{ye2026multistatebias}; and on GPCR peptide
complexes specifically, PAE over-estimates confidence precisely where the peptide has been placed
incorrectly, so confidence-first filtering cannot separate correct from incorrect placement
\citep{junker2026peptidedesign}. On the largest test of the question --- 92 fold-switching proteins
whose \emph{both} conformations were deposited before training --- the failure is not neutrality but
inversion: ``both plDDT and pTM scores assigned lower confidences to diverse correctly predicted
conformers and higher confidences to predictions that have not been observed experimentally'', and
the structure module ranks the lower-energy conformer above the higher-energy one ``50\% of the
time, equal to random chance'' \citep{chakravarty2024memorization}. Sampling harder and ranking by
confidence is therefore not merely uninformative but actively misleading, and this constrains
everything that follows: any method claiming a state must supply a state criterion independent of
the model that produced the structure.

Attempts to recover more than one state divide cleanly by where they intervene, and the oldest
family intervenes on the input. The simplest version samples harder without steering at all, either
by injecting dropout at inference and generating thousands of models per target
\citep{wallner2023afsample}, or by retraining a predictor on a conformationally split databank and
drawing on the order of a hundred samples per target, which recovers a held-out conformation for
52\% of targets but only as a best-of-N result with no rule for picking the right sample
\citep{bryant2024cfold}. The alignment has been manipulated in most of the ways available: reduced in
depth to as few as sixteen sequences, with recycling cut to one, which recovers both conformations of
five transporters and three GPCRs \citep{delalamo2022sampling}; clustered by
edit distance so that different subsets favour different conformations \citep{waymentsteele2024cluster},
a mechanism subsequently challenged as no better than random shallow subsampling
\citep{schafer2025confounds} and then defended with column-shuffling controls \citep{waymentsteele2025reply};
subsampled below the depth at which coevolutionary signal survives \citep{lee2025seqassoc}; edited by
in-silico alanine mutagenesis of alignment columns \citep{stein2022speachaf}; and masked at a tuned
column fraction \citep{kalakoti2025afsample2,kalakoti2026afsample3}, including masking restricted to the
neighbourhood of the orthosteric pocket in class A receptors \citep{mitjavila2026afsample2t}. Whether any of this reaches
the GPCR active state is contested, and the disagreement is instructive. Depth reduction recovers
both the inactive and the active conformation of three receptors spanning classes A, B1 and F
\citep{delalamo2022sampling}; tested head to head on $\beta_2$AR against AlphaFold3-lineage
backbones, however, ``MSA-level manipulation alone, whether through evolutionary clustering or
random subsampling, is largely insufficient to overcome the systematic conformational bias observed
across the deep learning tools evaluated in this study'' \citep[pp.~16--17]{ye2026multistatebias}.
Three differences separate them --- structural templates are supplied in the first and withheld in
the second, the backbone is AlphaFold2 rather than the AlphaFold3 lineage, and no receptor is shared
--- and none of the three has been varied deliberately. What the earlier work does settle is the
direction in which such a result must be read: its authors attribute the recovery to a training set
that ``featured the structures of many active GPCRs'' \citep{delalamo2022sampling}, so reaching the
active state is what the memorised prior already does, and only the inactive direction distinguishes
steering from recall.

The other input-side family supplies the answer rather than perturbing the question.
State-annotated GPCRdb templates combined with deletion of the alignment allow an operator to
obtain either activation state at near-experimental accuracy \citep{heo2022multistate}, and that route
has since been used to pin active-state templates for ranking peptide agonists and deorphanising
receptors \citep{ferguson2026deorphann}, and to condition a fine-tuned co-folding model on a declared
state for peptide design \citep{yang2025statespecific}. Its limits are documented by the people who use
it. Templates alone do not move a prediction when a deep alignment is present, because ``using such
curated template databases was not sufficient to make meaningful changes'' \citep[p.~1875]{heo2022multistate};
in transporters a correct alternative-state template is flipped back to the memorised state
\citep{swapna2025memorization}; and on a panel spanning kinases and class A receptors both handles
together fall short, where ``attempts to guide predictions through MSA editing or templating proved
to be insufficient, especially for the active-state conformation'' \citep[p.~18]{obendorf2026statespecific}.

A newer family leaves the inputs untouched and intervenes inside the generative process. Its
machinery comes from the diffusion literature rather than from structural biology: twisted
sequential Monte Carlo for asymptotically exact conditional sampling \citep{wu2023tds}, reward-weighted
particle steering that requires no differentiable objective \citep{singhal2025fksteering}, and
discriminator guidance in both the continuous \citep{kim2023refining} and the discrete
\citep{ekstromkelvinius2024discriminator} settings. None of those four papers contains a conformational
experiment and two contain no biomolecule at all, so they establish the mechanism rather than any
structural result. Conditioning a protein diffusion model on a classifier is itself not new: Chroma
samples proteins under external constraints including ``semantic specification from classifiers''
\citep[Discussion]{ingraham2023chroma}, though it designs new proteins rather than selecting among
conformations of an existing one. Applied to structure models, the same machinery takes several forms.
Interventions have been placed on the coordinates during denoising, by adding the gradient of a
differentiable collective variable to the sampler update \citep{lam2026metadiffusion}, or by running
particles under a bias away from the model's own default prediction and reweighting them
\citep{richman2025conformix}. Others act on the trunk representation, by perturbing the conditioning
tensors \citep{jung2026boltzperturb}, scaling the pair representation with a single scalar
\citep{suzuki2026pairscaling}, applying a mutual repulsion between parallel samples
\citep{suzuki2026conforflux}, or learning an affine transform of the pair representation
\citep{lee2026confornets}. Others still act on the conditioning embedding by gradient ascent
\citep{li2026embedding}, or on the alignment feature through a gradient taken back through the model's
own distogram \citep{tang2026steeraf}. Mechanistic work supports the choice of site, identifying the pair
track rather than the single track as the causal geometric substrate \citep{feldman2026alphainterp} and
showing that trunk representations are steerable across model families \citep{lu2026twostages}, with the
caveat that a concept being linearly decodable does not establish that the model uses it causally
\citep{jedryszek2026probing}. Sampling and steering machinery of this general kind now ships in
production co-folding code \citep{wohlwend2024boltz1,passaro2025boltz2}. A parallel strategy abandons
the predictor altogether and trains a generative emulator on molecular dynamics, which produces
quantitative equilibrium ensembles but was trained on soluble single chains rather than on membrane
proteins \citep{lewis2025bioemu}.

A smaller and, for the present purpose, more relevant body of work uses a biological co-input as
the handle. Supplying the cognate ligand of a cryptic site shifts AF3 ensembles onto the open
pocket while ligand-free runs return the closed one \citep{lazou2026cryptic}; co-folding a competitive
blocker restores discovery of allosteric and cryptic sites that the orthosteric pocket otherwise
overwrites \citep{purnomo2026cafe}; and across four multi-state proteins, including $\beta_2$AR,
``small-molecule ligands have weak or inconsistent effects, while large protein partners drive clear
conformational switching between states'' \citep[p.~2]{ye2026multistatebias}. The obvious objection to all of
these is that the co-input may be acting as an occupant rather than as itself, and one study has
tested it directly: across the same 82 enzymes, ligands known not to bind induce nearly the same
domain motion as the native trigger, and although they are placed with lower confidence, that
confidence is ``generally not sufficient for discriminating between binder and nonbinder ligands''
\citep[Significance]{yu2026domainmotion}. Any claim that a co-input sets a state therefore has to
separate the co-input's identity from its mere presence, and has to do so with a control arm rather
than an argument.

Running alongside all of this is a growing measurement literature: multi-state benchmarks \citep{ku2026promise,sun2026kinconfbench},
receptor-specific evaluations \citep{chib2025gpcrstates,zhang2026generalization,obendorf2026statespecific}, site-type comparisons holding the target fixed and varying only
whether the ligand is orthosteric or allosteric \citep{nittinger2025cofolding}, and explicit proposals
for the evidential bar that a multi-state claim should have to clear \citep{chakravarty2026statespace,liu2026ensembletests}.

Surveying this work by what it does rather than by what it studies exposes a consistent shape.
Prospectivity is close to absent. Of the 78 papers surveyed here, exactly three report an
unqualified prospective result, and none is of the relevant kind: a wet-lab peptide study that
generates no structures \citep{tran2026nanogs}, a blind CASP submission that defines no conformational
state at all \citep{wallner2023afsample}, and a de novo design campaign that generates new proteins rather
than alternative states of an existing one \citep{ingraham2023chroma}. Where a method does direct the state, the
direction is usually obtained from a structure of the answer. Cfold's recovery rate is a best-of-N
selected by TM-score against the held-out structure, and the paper demonstrates that no
confidence-based substitute exists \citep[p.~5]{bryant2024cfold}. ConforNets achieves transferable state
control, but the supervision signal for its transfer task is itself a deposited target structure
\citep{lee2026confornets}, so it cannot be pointed at a receptor whose target state has never been
solved. Trunk-scaling and repulsion strengths are swept on the same benchmarks that report the
results \citep{suzuki2026pairscaling,suzuki2026conforflux}. The cleanest exception is ConforMix, which
biases away from the model's own default prediction and so needs no reference structure at all, on
the explicit claim that prior conditional sampling on biomolecular diffusion models ``has required
additional input information, such as experimentally measured pairwise distances''
\citep[p.~2]{richman2025conformix}. Even there, coverage is scored against deposited references and both
reported rows are selected as the closest sample to one \citep[p.~6]{richman2025conformix}.

Where a directional handle exists at all, it is usually the operator's rather than the biology's.
Nineteen of the 78 papers carry one, and among prediction pipelines applied to receptors the state
is supplied by the operator in all but one case: as a state-annotated template \citep{heo2022multistate},
as a pinned active-state template \citep{ferguson2026deorphann}, as a declared state label routed through
state-matched templates \citep{yang2025statespecific}, or as a learned transform supervised by a
reference structure \citep{lee2026confornets}. The exception is treated below. Where the handle is
instead a biological co-input, one of three things is generally true. The system is not a receptor
\citep{lazou2026cryptic,purnomo2026cafe,richman2025conformix}; or the demonstration rests on four
targets with no memorization control \citep{ye2026multistatebias}; or the quantity measured is
binding-site geometry for docking rather than the receptor's activation state
\citep{mitjavila2026afsample2t}.

Two near-misses need stating precisely, because both are close and the gap has to be defined
against them rather than around them.

The first supplies the partner and measures the receptor, which is the combination this work is
built on. Comparing six AlphaFold protocols on 63 class A receptor--G$\mathrm{s}$ pairs released after the
models' training cutoffs, explicitly co-folding the receptor with its G protein reproduces the
intracellular rearrangement of activation better than operator-supplied active-state templates
with alignment masking or subsampling do, with ``TM5, TM6, and TM7 present[ing] three distinct
groups, with the models explicitly modeling G-protein binding outperforming the others''
\citep[p.~6302]{chiesa2025templatebias}. The advantage is not simply a training-set artefact: 94 of the
145 benchmark structures involve a receptor with no template and no presence in either model's
training set, and six receptor families are absent from both, which leads the authors to conclude
that the gain ``should not be attributed only to a more recent and active-state biased training
set, but rather to its ability in modeling allosteric effects during the cofolding process''
\citep[p.~6305]{chiesa2025templatebias}. A biological co-input has therefore already been shown to drive a
receptor's predicted activation state, and four narrower distinctions separate that demonstration
from the question asked here. The co-input there is the whole G$\alpha$ subunit or heterotrimer rather
than the isolated C-terminal segment, so the contribution of the $\alpha_5$ contact is never separated
from the rest of the interface. The state is scored as RMSD to the deposited active structure of
that same complex, so the procedure cannot be run on a receptor whose active state is unsolved.
There is no decoy or scrambled-partner arm, so occupancy and identity are not distinguished. And
no AlphaFold 3-lineage model is evaluated, a limitation the authors state themselves
\citep[p.~6299]{chiesa2025templatebias}.

The second near-miss approaches from the peptide side. A GPCR-specialised co-folding model
conditioned on activation state ranks designed peptides into agonists and antagonists, with
nanomolar hits assayed, and its authors identify state control as exactly what prior
peptide-design methods lack, writing that ``these methods do not allow for precise control over the
GPCR's functional state (e.g., active or inactive)'' \citep[p.~11426]{yang2025statespecific}. What separates
that work from the question asked here is the direction of the dependency. The state is declared
before any peptide is scored, and the authors are explicit that this is the operating mode:
``HF-Multistate, when specifying the GPCR state, can partially capture functional shifts in peptide
design'' \citep[p.~11434]{yang2025statespecific}. The peptide there is the designed output. No experiment in
that work varies a peptide sequence with the state conditioning held fixed and reports the
resulting receptor conformation.

That verification, however, does not scale. In the \emph{largest} GPCR and G protein co-folding studies
the receptor's state is not measured at all. Interface contact fingerprints have been extended to
825 modelled receptor and G protein complexes with no activation criterion defined anywhere in the
work \citep{matic2023gpcrome}; the entire GPCRome has been co-folded against all G$\alpha$ subunits with the
receptor's state assumed from the presence of the partner and no activation predicate applied
\citep{miglionico2026atlas}; and 5,595 AF2-Multimer complexes have been released gated purely on
self-confidence, with no activation state assigned to any model \citep{pandyszekeres2024gproteindb}. The
one study that does verify covers 63 pairs \citep{chiesa2025templatebias}, two orders of magnitude fewer
than the unverified resources. That asymmetry matters, because it is the large unverified
collections, not the small verified one, that downstream work treats as a source of active-state
structures. The instruments needed to close the loop already exist, in two
independent forms: a geometric index with published cut-offs that assigns class A activation state
from five inter-helical distances \citep{ibrahim2019a100}, and a learned classifier that separates
active from inactive deposited GPCR structures with near-perfect discrimination, whose authors
propose the missing experiment themselves --- that ``By applying HYALINE to AlphaFold 3-generated
GPCR models, researchers can obtain rapid, quantitative estimates of whether the predicted structure
represents an active, inactive, or intermediate conformation'' \citep[p.~11]{khaleq2026hyaline}. No
such experiment is reported there or anywhere else in this corpus, with either instrument.

Underlying all of this is a measurement problem that the field has not settled. Of the 31 papers
surveyed here that concern receptors, 13 either call the activation state from a structural overlay
by eye or define no state criterion at all, and the second group includes some of the largest
modelling efforts: 825 receptor--G protein complexes with no activation criterion defined anywhere
in the work \citep{matic2023gpcrome}, and 5{,}595 released complexes with no state assigned to any
model \citep{pandyszekeres2024gproteindb}. Among the studies that do quantify the transition, no
common admission rule has emerged. Reported TM6 displacements are anchored on different residues
and span an order of magnitude, from roughly 18\,\AA{} at the cytoplasmic end of TM6 in class B
GCGR \citep[p.~3]{hilger2020gcgr}, to about 14\,\AA{} for full activation of
$\beta_2$AR \citep[p.~13]{ye2026multistatebias}, to a combined TM6--TM7 shift of about
2\,\AA{} in a bistable opsin \citep{tejero2024opsin}. Each is a measurement of an observed
difference between two solved structures rather than a rule that could be applied to an unseen one.

Such a rule does exist, and its neglect is more telling than its absence would be. A universal
activation index for class A receptors, trained on microsecond simulations and validated against
268 crystal structures spanning 50 receptors, assigns state from five inter-helical
C$\alpha$--C$\alpha$ distances with published cut-offs, and its authors state that it ``can thus be
used to classify both experimental and computational GPCR structures''
\citep[p.~3938]{ibrahim2019a100}. It has been available since 2019, it is reimplementable from the
paper, and no study in this corpus applies it --- or any other predicate fixed in advance --- to a
predicted receptor structure. The gap is not that the field lacks an admission rule. It is that the
rules it has are used to describe structures that are already solved.

Some of that heterogeneity is principled rather than careless, and it constrains what a better
criterion can look like. Activation is not a switch: class A receptors are described as
interconverting among active, inactive and intermediate-active conformations on a rheostat
\citep[p.~3696]{georgiou2025heterogeneity}, so a single hard cut-off imposed on a continuum can
destroy the signal it is meant to detect. A usable state criterion therefore has to be fixed before
the data are seen, reported with its threshold, and checked against a second, independent readout
so that its disagreement rate is known rather than assumed. No study in this corpus reports such a
rate.

The benchmarks that might otherwise have settled the question exclude receptor activation by
construction, and they now do so three times over by three different mechanisms. ProMiSE states it
in its own Limitations: ``Our pair-extraction criterion, TM-score $<$ 0.8, is also stringent and
preferentially captures large conformational changes. Consequently, localized but biologically
important motions, such as GPCR TM6 displacement or transporter pocket rearrangements, may be
excluded'' \citep[p.~9]{ku2026promise}. The same TM-score 0.8 rule defines the structural clusters of the
Cfold benchmark \citep[p.~7]{bryant2024cfold} and is inherited without comment by work built on it
\citep{kalakoti2026afsample3}, whose two-state targets are taken from that benchmark's own release. The
strongest inference-time sampling result to date arrives at the same place by a third route,
importing its four evaluation sets wholesale from earlier benchmark papers, none of which contains
a receptor \citep{richman2025conformix}. ProMiSE additionally reports that partners are the harder case:
holo-conditioned prediction fails for 4 to 17\% of pairs when the partner is a small molecule but
for 34 to 57\% when it is a protein chain, and the two failure modes differ in kind, since
``ligand-induced failures are \emph{asymmetric} (collapse in a single direction, Apo$\rightarrow$Holo), whereas
protein-induced failures are \emph{symmetric}, with substantial collapse in both directions''
\citep[p.~17]{ku2026promise}. That result runs opposite to the axis examined here and is addressed directly
in the Discussion.

The last column separates two questions that are usually merged, and it is where the receptor
literature is weakest. Forty-four of the 78 papers have some held-out or post-cutoff set in their
design; far fewer run that set as an analysed control arm, and among the papers that actually
address GPCR conformational state the control is absent, unpowered or unmatched
\citep{chib2025gpcrstates,ye2026multistatebias,obendorf2026statespecific,yang2025statespecific}. The
contrast with what a properly controlled design looks like is stark: the enzyme study above
stratifies its 82 targets by the composition of the training set itself, runs a non-binding-ligand
arm against the native trigger, and replicates the whole result in a second architecture
\citep[Results]{yu2026domainmotion}. Even
the cleanest inference-time result concedes the point, noting that its evaluated proteins ``were
likely present in the Boltz training set'' \citep[p.~6]{richman2025conformix}. The omission is consequential
rather than pedantic. Co-folding's advantage over physics-based methods is concentrated near the
training distribution and largely disappears in the least-similar stratum
\citep{skrinjar2026generalization,roehrig2026docking}; splitting on sequence identity does not by itself
stop benchmark leakage \citep{mattsson2026leakage}; and predicted ligand placement can survive a pocket
being mutated until it cannot bind, which shows retrieval rather than physics driving the result
\citep{masters2025physics}. That last failure extends beyond the structural head: on 943
virtual-screening hits across ten targets, most of them receptors, Boltz-2's separation of active from
inactive compounds ``remains insensitive to key binding site mutations and even in some cases to
target exchange'' \citep[Abstract]{bret2025boltz2docking}. A conformational claim carrying no matched
control arm is therefore not yet distinguishable from a lookup, and a control arm is only
persuasive if the quantity it perturbs is one the model demonstrably reads.

Taken together, the gap is narrower than a survey of this literature first suggests, and it is
sharper for being narrow. A biological co-input has been shown to set a receptor's predicted
conformational state, and that state has been verified, but only with the whole transducer
supplied and only by measuring distance to the deposited answer. What the corpus does not contain
is a study in which the co-input is reduced to the minimal contact that carries activation, in
which its contribution is separated from mere occupancy of the cleft by a graded series of decoy
and sequence-scrambled controls, and in which the resulting state is called by a predicate fixed
in advance rather than by proximity to a structure that must already exist.

That gap has three parts and this work closes two of them. The graded decoy and sequence-scrambled
series is run, and the state is called by a predicate fixed in advance and cross-checked against a
second instrument. The third part -- reducing the co-input to the minimal contact that carries
activation -- is motivated here but \emph{not} delivered: every arm reported below supplies a whole
G$\alpha$ subunit or a nanobody, and no arm supplies a peptide of any length.

\textbf{[PI: this is the decision that sets the title, and it has two parts. First, the segment
actually manipulated is \emph{eleven} residues, not twenty-one --- the decoy tail is eleven, the
drop's varied quantity is the eleven-residue tail helicity, and 4X1H, the only peptide-bound entry
in the reference set, is an engineered eleven-mer. Second, no arm supplies a peptide at all, so a
narrowed title is not about a peptide of any length: it is about a G$\alpha$ co-input and its
C-terminal determinant. That is a smaller claim than the current title and a much better supported
one. The alternative is to run a peptide arm. See \texttt{../CLAIMS.md} C6 and C7.]}


\subsection{What we do}

\begin{quote}
\textbf{PROVISIONAL STUB.} Landed blocks only, no numbers, and expected to be rewritten once the data
export is complete. See \texttt{../STATUS.md} for what has landed and what has not.
\end{quote}

We treat a G$\alpha$ transducer supplied alongside the receptor as a co-input to a frozen co-folding
model and ask what the receptor does in response, holding the model weights, the alignment and the
template settings unchanged and varying only what is supplied alongside the receptor sequence.
\textbf{What is supplied is the whole G$\alpha$ subunit, or a nanobody, not the 21-residue
$\alpha_5$ C-terminal peptide of the title; no arm in any landed block supplies a peptide.}
Predictions are run across a panel of class A receptors on four independent co-folding backbones,
and the receptor's conformational state is called by a predicate defined in advance rather than by
eye. Four blocks have landed. The first compares apo receptors against receptors given their cognate
partner. The second places apo, decoy, sequence-shuffled and cognate co-inputs on a common scale. The decoy
shares the cognate subunit's entire scaffold and differs from it only in the eleven C-terminal
residues, which are permuted rather than replaced, so composition is held constant and only the tail
moves. That construction gives a decomposition in two steps that telescope exactly: supplying the
subunit at all, with its C-terminal sequence scrambled, and then restoring that sequence. The second
step is the contribution of the $\alpha_5$ C terminus itself, isolated against an otherwise
identical partner. Whether the sequence must additionally belong to the \emph{correct} G$\alpha$
family is a separate and weaker question, and the shuffled arm is where it is asked. The third asks whether the orthosteric pocket, as distinct
from the cytoplasmic face, separates by ligand class. The fourth varies the depth of the sequence
alignment itself, asking whether starving the model of homologue data moves the same predicate --
and, on a second and finer measure, whether a predicate that clears under shallow alignments
corresponds to a structure that has actually reached the target state. Alongside these we report a census of
intermediate conformations produced under partner variation, and an observation about model
confidence made across backbones on a single receptor.

\subsection{Contributions}

\begin{quote}
\textbf{PROVISIONAL STUB.} Four candidate contributions, each written so that it can be falsified.
The wording will change with the final result set; the framing is intended to survive it.
\end{quote}

\begin{enumerate}
  \item A biological co-input, supplied the way a user of these models would supply it, moves the predicted conformational state of a receptor. Falsified if the state distribution under the cognate co-input is indistinguishable from the apo distribution.
  \item The effect is attributable to the co-input rather than to steric occupancy alone, because the decoy and sequence-shuffled arms sit between the two extremes. Falsified if a decoy of matched size reproduces the cognate effect.
  \item The state is verified rather than assumed, by a predicate defined in advance and reported with its threshold, applied identically to every arm and cross-checked against a second instrument. Falsified if the two instruments disagree on which arms the co-input separates --- that is, on the ordering of apo, decoy and the partner-bearing arms. The two rank \emph{shuffled} against \emph{cognate} differently, and that is not a falsification: those two arms sit about 0.1~\AA{} apart on the tilt axis, which is inside the instrument's resolution, and both instruments agree on apo $<$ decoy $<$ \{shuffled, cognate\} on every backbone.
  \item Model confidence does not track conformational correctness on this axis, so it cannot be used to select a state. Falsified if confidence separates correct from incorrect state calls at matched partner composition.
\end{enumerate}